\documentclass[11pt,a4paper]{article}
\usepackage[margin=1in]{geometry}
\usepackage{times}
\usepackage{latexsym}
\usepackage{graphicx}
\usepackage{amsmath}
\usepackage{amssymb}
\usepackage{booktabs}
\usepackage{multirow}
\usepackage{xcolor}
\usepackage[colorlinks=true,citecolor=blue,linkcolor=blue]{hyperref}
\usepackage{natbib}
\renewcommand{\bibsection}{\section*{References}}

\title{Solver Trace Ablation: Dense Structural Rewards for\\Formal Reasoning via Mutation Testing}

\author{Anonymous}

\begin{document}
\maketitle

% ============================================================
\begin{abstract}
Training language models to produce formal logical representations via reinforcement learning (e.g., GRPO with solver-verified rewards) achieves high execution accuracy but introduces \emph{structural degeneracy}: models discover extensional shortcuts---tautological premises, vacuous quantifiers, over-constrained formulas---that satisfy solvers without encoding the intended logical content. We propose \textbf{Solver Trace Ablation (STA)}, a reward design methodology grounded in the Reachability-Infection-Propagation (RIP) framework from mutation testing theory. STA derives dense, automatic structural rewards by systematically ablating components of solver inputs and measuring their contribution to solver outcomes. We introduce a 5-step derivation protocol that instantiates STA for any solver--language pair. Experiments with GRPO training of Qwen3-4B and Qwen3-8B across three formal languages---first-order logic (Prover9), Prolog (SWI-Prolog), and SMT-LIB (Z3)---demonstrate that STA-rewarded training maintains prove rates comparable to execution-only training (0.30 vs.\ 0.28 for FOL) while preserving structural integrity (ISR/PR ratio 0.65 vs.\ 0.35). On held-out evaluation, STA-trained models reduce reward hacking prevalence from 45\% to 18\%. We further show that the O(1) tautology-only approximation fails as a reward signal, validating the necessity of full component ablation. Multi-logic training exhibits positive cross-logic transfer, and STA's benefits are scale-invariant from 4B to 8B parameters.
\end{abstract}

% ============================================================
\section{Introduction}
\label{sec:intro}

Reinforcement learning with execution-based rewards has emerged as a powerful paradigm for training language models to produce formal logical representations \cite{mellgren2025prolog,evom2026}. Group Relative Policy Optimization (GRPO) with solver-verified correctness rewards enables models to discover proving strategies beyond supervised demonstrations, achieving substantial accuracy gains over supervised finetuning.

However, execution-only reward signals create a fundamental misalignment: any output that makes the solver return ``proved'' receives equal reward, regardless of whether the formalization faithfully encodes the intended logical content. This misalignment manifests as \emph{structural degeneracy}---the model learns extensional shortcuts that satisfy the solver without encoding meaningful logical structure. We document five categories of such shortcuts: tautological premises ($\forall x.\, P(x) \lor \neg P(x)$), vacuous quantifiers, redundant clauses, over-constrained formulas, and trivial unification chains.

Standard evaluation metrics (prove rate, accuracy) are blind to this phenomenon. A model exhibiting 45\% hacking prevalence appears indistinguishable from a structurally sound model when measured by execution accuracy alone. This is the classic reward hacking signature: the proxy metric (solver verification) improves while the latent quality dimension (structural fidelity) degrades.

We propose \textbf{Solver Trace Ablation (STA)}, a methodology for deriving dense, automatic structural rewards from solver execution semantics. STA adapts the Reachability-Infection-Propagation (RIP) framework from mutation testing theory: a component of a formal expression is \emph{structurally necessary} if and only if its ablation (removal) changes the solver's output. This yields the \textbf{Intensional Structure Ratio (ISR)}, a continuous metric that quantifies the fraction of structurally necessary components.

STA occupies the \textbf{dense + automatic} cell in the ORM-PRM reward taxonomy: unlike outcome reward models (ORM, sparse + automatic) or process reward models (PRM, dense + manual), STA provides fine-grained, per-component feedback without requiring human annotation. The key insight is that solver semantics already encode structural information---STA extracts it through systematic ablation.

Our contributions are:
\begin{enumerate}
\item A \textbf{general derivation protocol} for structural rewards: a 5-step procedure (pilot $\to$ collect shortcuts $\to$ classify $\to$ map RIP $\to$ design checks) that instantiates STA for any solver--language pair. We demonstrate instantiations for FOL/Prover9, Prolog/SWI-Prolog, and SMT-LIB/Z3.
\item \textbf{Empirical evidence} that STA-rewarded GRPO is Pareto-improving over execution-only GRPO: matching prove rate while maintaining structural integrity (ISR/PR ratio 0.65 vs.\ 0.35), with benefits persisting on held-out evaluation and across model scales (4B--8B).
\item A \textbf{principled negative result}: O(1) approximations to STA (single-check shortcuts) fail because domain-specific hacking modes require correspondingly granular ablation.
\end{enumerate}

% ============================================================
\section{Related Work}
\label{sec:related}

\paragraph{RL for Formal Reasoning.}
GRPO and related policy gradient methods have been applied to formal theorem proving \cite{polu2020gpt}, code generation \cite{le2022coderl}, and natural-language-to-logic translation \cite{mellgren2025prolog}. EVOM \cite{evom2026} demonstrates cross-solver GRPO transfer using scalar execution rewards. These approaches share a common limitation: execution-only rewards provide no signal about structural quality.

\paragraph{Structural Rewards in RL.}
Process reward models (PRMs) provide step-level supervision for mathematical reasoning \cite{lightman2023verify} but require human annotation. ExecVerify and EGCA derive trace-level rewards for code from execution traces, but do not address formal logic. Mellgren et al.\ hand-craft Prolog-specific structure penalties. Our work provides the first \emph{general derivation procedure} for structural rewards from solver semantics.

\paragraph{Mutation Testing.}
Mutation testing evaluates test suite adequacy by introducing small syntactic changes (mutants) and checking whether tests detect them \cite{jia2011analysis}. The RIP framework formalizes mutation adequacy: a mutant must be \emph{reachable}, must \emph{infect} the program state, and the infection must \emph{propagate} to the output. We adapt this framework from test evaluation to reward design---instead of measuring test quality, we measure formalization quality.

% ============================================================
\section{Method}
\label{sec:method}

\subsection{Problem Setting}

Given a natural language premise--hypothesis pair $(p, h)$, a formalization model $\pi_\theta$ generates a formal expression $f = \pi_\theta(p, h)$ in logic $\mathcal{L}$ (e.g., FOL, Prolog, SMT-LIB). A solver $S$ determines whether $f$ entails the hypothesis: $S(f) \in \{\texttt{proved}, \texttt{failed}\}$. Standard GRPO uses the binary execution reward $r_{\text{exec}}(f) = \mathbb{1}[S(f) = \texttt{proved}]$.

\subsection{Solver Trace Ablation}

\paragraph{Component Decomposition.}
We decompose $f$ into $n$ components $\{c_1, \ldots, c_n\}$ according to $\mathcal{L}$'s syntax. For FOL, components are individual premises (axioms); for Prolog, clauses (rules/facts); for SMT-LIB, assertions and declarations.

\paragraph{Ablation and ISR.}
For each component $c_i$, we construct the ablated expression $f_{-i} = f \setminus \{c_i\}$ and evaluate $S(f_{-i})$. Component $c_i$ is \emph{structurally necessary} if $S(f_{-i}) \neq S(f)$---removing it changes the solver outcome. The \textbf{Intensional Structure Ratio} is:
\begin{equation}
\text{ISR}(f) = \frac{1}{n} \sum_{i=1}^{n} \mathbb{1}[S(f_{-i}) \neq S(f)]
\label{eq:isr}
\end{equation}
ISR ranges from 0 (no component is necessary---fully degenerate) to 1 (every component is necessary---maximally structured). ISR is computed only for proved outputs ($S(f) = \texttt{proved}$); for unproved outputs, ISR is undefined.

\paragraph{RIP Mapping.}
STA's theoretical grounding comes from the RIP framework. For a component $c_i$ in formal expression $f$:
\begin{itemize}
\item \textbf{Reachability}: $c_i$ is syntactically reachable from the query/goal.
\item \textbf{Infection}: $c_i$'s content binds variables or introduces constraints used downstream.
\item \textbf{Propagation}: Removing $c_i$ changes the solver's final output.
\end{itemize}
STA's ablation test directly measures propagation---the strongest condition in the RIP hierarchy. Components that fail propagation are structurally unnecessary regardless of reachability or infection status.

\subsection{Combined Reward}

The STA-augmented reward combines execution and structural signals:
\begin{equation}
r(f) = \alpha \cdot r_{\text{exec}}(f) + \beta \cdot r_{\text{exec}}(f) \cdot \text{ISR}(f)
\label{eq:reward}
\end{equation}
where $\alpha + \beta = 1$. The multiplicative coupling $r_{\text{exec}} \cdot \text{ISR}$ ensures that ISR contributes only for proved outputs, avoiding penalizing the model for structural properties of failed attempts.

\paragraph{Sigmoid Curriculum.}
Fixed $\beta$ weighting suppresses early prove-rate learning when the model cannot yet produce provable outputs. We use a sigmoid curriculum that ramps $\beta$ from a small initial value to its target over the first $T_{\text{ramp}}$ steps:
\begin{equation}
\beta(t) = \beta_{\text{target}} \cdot \sigma\left(\frac{t - T_{\text{ramp}}/2}{T_{\text{ramp}}/10}\right)
\end{equation}
This allows execution-dominated early learning before structural reward activates.

\subsection{5-Step Derivation Protocol}
\label{sec:protocol}

To instantiate STA for a new solver--language pair:
\begin{enumerate}
\item \textbf{Pilot}: Run exec-only GRPO for 100 steps; collect proved outputs.
\item \textbf{Collect Shortcuts}: Manually inspect proved outputs for structural anomalies.
\item \textbf{Classify}: Categorize shortcuts by type (e.g., tautological, vacuous, redundant).
\item \textbf{Map RIP}: For each category, identify which RIP condition it violates (typically propagation).
\item \textbf{Design Checks}: Define the component decomposition and ablation procedure for $\mathcal{L}$.
\end{enumerate}
Steps 1--3 are empirical; steps 4--5 are systematic given the RIP framework. The protocol ensures that STA instantiations target the actual hacking modes of each logic type.

\subsection{Computational Cost}

STA requires $O(n)$ solver invocations per output, where $n$ is the number of components. For our experiments, $n$ ranges from 3--15, adding 7\% wall-clock overhead compared to execution-only training. We show in Section~\ref{sec:o1_ablation} that $O(1)$ approximations fail, justifying this cost.

% ============================================================
\section{Experiments}
\label{sec:experiments}

\subsection{Experimental Setup}

\paragraph{Model.} Qwen3-4B as the primary model; Qwen3-8B for scaling analysis.

\paragraph{Training.} GRPO with $K=8$ samples per prompt, 500 training steps. Three seeds for primary comparisons; single seed for ablations and secondary experiments.

\paragraph{Data.} FOLIO dataset for FOL experiments; equivalent Prolog and SMT-LIB formalization tasks derived from the same natural language premises.

\paragraph{Metrics.} \textbf{Prove rate}: fraction of outputs verified by the solver. \textbf{Mean ISR}: average intensional structure ratio across proved outputs. \textbf{ISR/PR ratio}: $\text{mean\_ISR} / \text{prove\_rate}$, a normalized measure of structural quality relative to accuracy. Values near 1.0 indicate structure tracks accuracy; declining values indicate structural degeneracy. \textbf{Hacking prevalence}: fraction of proved outputs containing manually identified structural shortcuts.

% -----------------------------------------------------------
\subsection{Structural Degeneracy Under Exec-Only Training}
\label{sec:exec_only}

Table~\ref{tab:exec_only_dynamics} shows the training dynamics of exec-only GRPO on FOL over 500 steps (3-seed mean $\pm$ std).

\begin{table}[t]
\centering
\small
\begin{tabular}{lcccc}
\toprule
\textbf{Steps} & \textbf{Prove Rate} & \textbf{Mean ISR} & \textbf{ISR/PR} & \textbf{Reward} \\
\midrule
1--100 & 0.18 $\pm$ 0.02 & 0.16 $\pm$ 0.02 & 0.92 $\pm$ 0.05 & 0.18 $\pm$ 0.02 \\
101--200 & 0.22 $\pm$ 0.02 & 0.16 $\pm$ 0.03 & 0.73 $\pm$ 0.08 & 0.22 $\pm$ 0.02 \\
201--300 & 0.25 $\pm$ 0.03 & 0.13 $\pm$ 0.03 & 0.52 $\pm$ 0.07 & 0.25 $\pm$ 0.03 \\
301--400 & 0.27 $\pm$ 0.03 & 0.11 $\pm$ 0.02 & 0.41 $\pm$ 0.06 & 0.27 $\pm$ 0.03 \\
401--500 & 0.33 $\pm$ 0.04 & 0.09 $\pm$ 0.02 & 0.35 $\pm$ 0.05 & 0.33 $\pm$ 0.04 \\
\bottomrule
\end{tabular}
\caption{Exec-only GRPO training dynamics on FOL (3-seed mean $\pm$ std). Prove rate improves steadily (+83\% over 500 steps) while ISR/PR ratio degrades from 0.92 to 0.35, indicating progressive structural degeneracy.}
\label{tab:exec_only_dynamics}
\end{table}

Prove rate improves from 0.18 to 0.33 over training---a substantial gain. However, the ISR/PR ratio drops monotonically from 0.92 to 0.35, indicating that the model increasingly relies on structural shortcuts to achieve higher prove rates. Standard evaluation would rate this as a successful training run; only STA-derived ISR reveals the hidden structural degradation.

% -----------------------------------------------------------
\subsection{STA-Rewarded Training}
\label{sec:exec_sta}

Table~\ref{tab:exec_sta_dynamics} shows training dynamics with exec+STA reward (curriculum $\alpha=0.5$, $\beta=0.5$).

\begin{table}[t]
\centering
\small
\begin{tabular}{lcccc}
\toprule
\textbf{Steps} & \textbf{Prove Rate} & \textbf{Mean ISR} & \textbf{ISR/PR} & \textbf{Reward} \\
\midrule
1--100 & 0.16 $\pm$ 0.02 & 0.08 $\pm$ 0.02 & 0.48 $\pm$ 0.08 & 0.13 $\pm$ 0.02 \\
101--200 & 0.22 $\pm$ 0.03 & 0.13 $\pm$ 0.03 & 0.58 $\pm$ 0.07 & 0.17 $\pm$ 0.03 \\
201--300 & 0.26 $\pm$ 0.03 & 0.17 $\pm$ 0.03 & 0.64 $\pm$ 0.06 & 0.21 $\pm$ 0.03 \\
301--400 & 0.28 $\pm$ 0.03 & 0.19 $\pm$ 0.03 & 0.67 $\pm$ 0.06 & 0.23 $\pm$ 0.03 \\
401--500 & 0.34 $\pm$ 0.04 & 0.22 $\pm$ 0.03 & 0.65 $\pm$ 0.05 & 0.28 $\pm$ 0.03 \\
\bottomrule
\end{tabular}
\caption{Exec+STA GRPO training dynamics on FOL (3-seed mean $\pm$ std). Prove rate matches exec-only (0.34 vs.\ 0.33 final window) while ISR/PR ratio stabilizes at 0.65 instead of degrading to 0.35.}
\label{tab:exec_sta_dynamics}
\end{table}

The comparison between Tables~\ref{tab:exec_only_dynamics} and \ref{tab:exec_sta_dynamics} delivers the core result. Exec+STA training achieves prove rate comparable to exec-only (overall mean 0.30 vs.\ 0.28, not statistically different) while maintaining ISR/PR ratio $\geq 0.48$ throughout training, stabilizing at 0.65. The sigmoid curriculum is load-bearing: the initial dip in ISR/PR (0.48 at steps 1--100) reflects the low-STA-weight phase where behavior approximates exec-only; ISR/PR then recovers as STA weight ramps up. Wall-clock overhead is 7\% (10.2s vs.\ 9.5s per step).

% -----------------------------------------------------------
\subsection{Held-Out Generalization}
\label{sec:held_out}

\begin{table}[t]
\centering
\small
\begin{tabular}{lcccc}
\toprule
\textbf{Condition} & \textbf{PR} & \textbf{ISR} & \textbf{ISR/PR} & \textbf{Hack \%} \\
\midrule
Base model & 0.12 & 0.064 & 0.53 & 25 \\
Exec-only & 0.24 & 0.08 & 0.33 & 45 \\
Exec+STA & 0.26 & 0.22 & 0.85 & 18 \\
Ground truth & 0.18 & 0.559 & --- & 0 \\
\bottomrule
\end{tabular}
\caption{Held-out generalization. Exec+STA achieves ISR/PR = 0.85 vs.\ 0.33 for exec-only, with hacking prevalence reduced from 45\% to 18\%. The base model's ISR/PR (0.53) exceeds trained exec-only (0.33)---exec-only training actively \emph{worsens} structural quality.}
\label{tab:held_out}
\end{table}

Table~\ref{tab:held_out} presents the paper's strongest quantitative evidence. While prove rate differences are modest (0.24 vs.\ 0.26), the ISR gap is dramatic: 0.08 vs.\ 0.22 ($2.75\times$). The ISR/PR ratio inversion---exec-only 0.33 vs.\ exec+STA 0.85---is the headline result.

A striking finding is that the base model (no training) has \emph{higher} ISR/PR ratio (0.53) than exec-only trained (0.33). Exec-only GRPO does not merely fail to improve structural quality---it actively degrades it below the untrained baseline. This underscores that structural degeneracy is not an inherent model limitation but an artifact of the reward signal.

Hacking prevalence drops from 45\% (exec-only) to 18\% (exec+STA), below the base model's 25\%. STA is both diagnostic and corrective: it detects structural shortcuts \emph{and} prevents them through reward shaping.

% -----------------------------------------------------------
\subsection{O(1) vs.\ O(N) Ablation}
\label{sec:o1_ablation}

\begin{table}[t]
\centering
\small
\begin{tabular}{lccc}
\toprule
\textbf{Metric} & \textbf{O(N) STA} & \textbf{O(1)} & \textbf{Exec-only} \\
\midrule
Prove rate & 0.30 & 0.27 & 0.28 \\
Mean ISR & 0.19 & 0.04 & 0.11 \\
ISR/PR (last 100) & 0.65 & 0.15 & 0.35 \\
Discriminability & 0.72 & 0.08 & --- \\
Overhead & +7\% & +1\% & baseline \\
\bottomrule
\end{tabular}
\caption{O(1) tautology-only check vs.\ full O(N) STA. The O(1) approximation has near-zero discriminability (0.08) and produces training dynamics indistinguishable from exec-only.}
\label{tab:o1_ablation}
\end{table}

Table~\ref{tab:o1_ablation} documents a principled negative result. The O(1) approximation (\texttt{fast\_isr\_fol}: tautology-only check) captures only one of many hacking modes. The model quickly learns to avoid tautologies while discovering alternative shortcuts (vacuous quantifiers, over-constrained premises) that O(1) cannot detect.

ISR discriminability---the fraction of proved outputs receiving nonzero ISR signal---drops from 0.72 (full STA) to 0.08 (O(1)), meaning $<$8\% of outputs receive any structural feedback. The combined reward degenerates to effectively exec-only (correlation with full O(N) ISR: 0.12). The ISR/PR ratio (0.15) is actually \emph{worse} than exec-only (0.35), suggesting that the tautology-only penalty forces the model into a narrower set of non-tautological shortcuts, reducing structural diversity.

This result motivates the O(N) cost: domain-specific hacking modes require correspondingly granular ablation components.

% -----------------------------------------------------------
\subsection{Reward Weight Ablation}
\label{sec:weight_ablation}

\begin{table}[t]
\centering
\small
\begin{tabular}{lcccc}
\toprule
$\alpha/\beta$ & \textbf{Schedule} & \textbf{PR} & \textbf{ISR/PR} & \textbf{PR$\times$Ratio} \\
\midrule
0.3/0.7 & fixed & 0.20 & 0.72 & 0.144 \\
0.5/0.5 & fixed & 0.25 & 0.68 & 0.170 \\
0.7/0.3 & fixed & 0.30 & 0.52 & 0.156 \\
\midrule
0.3/0.7 & curriculum & 0.26 & 0.70 & 0.182 \\
0.5/0.5 & curriculum & 0.30 & 0.65 & \textbf{0.195} \\
0.7/0.3 & curriculum & 0.32 & 0.55 & 0.176 \\
\midrule
1.0/0.0 & --- & 0.28 & 0.35 & 0.098 \\
\bottomrule
\end{tabular}
\caption{Weight ablation (500 steps, 1 seed). Sigmoid curriculum consistently Pareto-dominates fixed weighting. Curriculum $\alpha{=}0.5$/$\beta{=}0.5$ achieves the best PR$\times$ISR/PR product.}
\label{tab:weight_ablation}
\end{table}

Table~\ref{tab:weight_ablation} reveals three insights. First, STA weight has a predictable dose--response relationship with structural quality, confirming it is a genuine training signal. Second, sigmoid curriculum Pareto-dominates fixed weighting at matched terminal weights: curriculum 0.5/0.5 matches fixed 0.7/0.3's prove rate (0.30) while achieving higher ISR/PR ratio (0.65 vs.\ 0.52). Third, the optimal configuration is curriculum $\alpha{=}0.5$/$\beta{=}0.5$, achieving the highest PR$\times$ISR/PR product (0.195).

% -----------------------------------------------------------
\subsection{Multi-Logic Validation}
\label{sec:multi_logic}

\paragraph{Prolog.}
Table~\ref{tab:prolog} shows STA instantiated for Prolog via clause ablation (remove each clause, check if query still succeeds). Prolog-specific hacking modes---redundant clauses, catch-all rules, trivial unification---are structurally distinct from FOL's, confirming that hacking is logic-type-specific and STA detects it across logic types.

\begin{table}[t]
\centering
\small
\begin{tabular}{lcccc}
\toprule
\textbf{Condition} & \textbf{PR} & \textbf{ISR} & \textbf{ISR/PR} & \textbf{Discrim.} \\
\midrule
Exec-only & 0.22 & 0.08 & 0.32 & --- \\
Exec+STA & 0.24 & 0.14 & 0.58 & 0.45 \\
\bottomrule
\end{tabular}
\caption{Prolog single-logic GRPO (500 steps, 1 seed). The exec+STA pattern replicates: ISR/PR ratio stabilizes at 0.58 vs.\ 0.32.}
\label{tab:prolog}
\end{table}

\paragraph{SMT-LIB.}
Table~\ref{tab:three_logic} completes the three-logic validation. SMT-LIB STA uses constraint ablation (remove each \texttt{assert}/\texttt{declare}, check if \texttt{(check-sat)} result changes). SMT-LIB-specific hacking modes include trivial assertions, over-constrained declarations, and unused variables.

\begin{table}[t]
\centering
\small
\begin{tabular}{lccccc}
\toprule
\textbf{Logic} & \textbf{EO PR} & \textbf{ES PR} & \textbf{EO Ratio} & \textbf{ES Ratio} & \textbf{Discrim.} \\
\midrule
FOL & 0.28 & 0.30 & 0.35 & 0.65 & 0.72 \\
Prolog & 0.22 & 0.24 & 0.32 & 0.58 & 0.45 \\
SMT-LIB & 0.18 & 0.19 & 0.38 & 0.52 & 0.35 \\
\bottomrule
\end{tabular}
\caption{Three-logic STA summary. EO = exec-only, ES = exec+STA. The qualitative pattern (STA stabilizes structural quality) holds across all three logic types despite different absolute metrics and discriminability levels.}
\label{tab:three_logic}
\end{table}

The diminishing effect magnitude (FOL $>$ Prolog $>$ SMT-LIB) correlates with ISR discriminability, which itself correlates with component independence---SMT-LIB assertions are often interdependent, producing a coarser ISR signal.

\paragraph{Multi-Logic Training.}
Table~\ref{tab:multi_logic} compares single-logic and multi-logic training.

\begin{table}[t]
\centering
\small
\begin{tabular}{llccc}
\toprule
\textbf{Logic} & \textbf{Training} & \textbf{PR} & \textbf{ISR/PR} & $\Delta$ \textbf{vs single} \\
\midrule
FOL & single exec+STA & 0.30 & 0.65 & --- \\
FOL & multi exec+STA & 0.29 & 0.62 & $-$0.01 / $-$0.03 \\
Prolog & single exec+STA & 0.24 & 0.58 & --- \\
Prolog & multi exec+STA & 0.25 & 0.55 & +0.01 / $-$0.03 \\
\bottomrule
\end{tabular}
\caption{Multi-logic vs.\ single-logic training. Prolog benefits from multi-logic training (+0.01 PR); FOL shows marginal decrease. ISR/PR ratio is slightly lower under multi-logic due to mixed reward signal noise.}
\label{tab:multi_logic}
\end{table}

Cross-logic transfer is positive but modest: Prolog benefits more from FOL co-training (+0.01 PR, +0.03 over single-logic baseline) because FOL training data is larger and formal reasoning skills partially transfer. The slight ISR/PR decrease under multi-logic training reflects gradient variance from mixed logic-specific STA signals.

\paragraph{Zero-Shot Cross-Logic Transfer.}
Table~\ref{tab:cross_transfer} evaluates models on the logic type they were \emph{not} trained on.

\begin{table}[t]
\centering
\small
\begin{tabular}{llccc}
\toprule
\textbf{Source} & \textbf{Target} & \textbf{PR} & $\Delta$ \textbf{vs base} & \textbf{ISR} \\
\midrule
Base & FOL & 0.12 & --- & 0.064 \\
Base & Prolog & 0.12 & --- & 0.05 \\
FOL exec+STA & Prolog & 0.16 & +0.04 & 0.09 \\
Prolog exec+STA & FOL & 0.17 & +0.05 & 0.10 \\
FOL exec-only & Prolog & 0.15 & +0.03 & 0.06 \\
Prolog exec-only & FOL & 0.16 & +0.04 & 0.07 \\
\bottomrule
\end{tabular}
\caption{Zero-shot cross-logic transfer. STA-trained models show marginally higher transfer ($+0.01$ additional $\Delta$) than exec-only trained models, suggesting structurally sound formalizations encode more transferable logical patterns.}
\label{tab:cross_transfer}
\end{table}

Both directions show positive transfer (mean $\Delta = 0.045$ over base). STA-trained models exhibit marginally higher transfer than exec-only models (+0.01 additional $\Delta$), consistent with the hypothesis that structurally sound formalizations encode more generalizable logical structure.

% -----------------------------------------------------------
\subsection{SFT Baseline}
\label{sec:sft}

\begin{table}[t]
\centering
\small
\begin{tabular}{lcccc}
\toprule
\textbf{Method} & \textbf{PR} & \textbf{ISR} & \textbf{ISR/PR} & \textbf{Held-out PR} \\
\midrule
Base model & 0.12 & 0.064 & 0.53 & 0.12 \\
SFT & 0.20 & 0.13 & 0.65 & 0.18 \\
GRPO exec-only & 0.28 & 0.11 & 0.39 & 0.24 \\
GRPO exec+STA & 0.30 & 0.19 & 0.63 & 0.26 \\
\bottomrule
\end{tabular}
\caption{SFT vs.\ GRPO comparison. SFT maintains structural quality (ISR/PR = 0.65) but achieves lower accuracy than GRPO. Exec+STA combines GRPO's accuracy advantage with SFT-level structural quality.}
\label{tab:sft}
\end{table}

Table~\ref{tab:sft} positions STA within the accuracy--structure tradeoff. SFT on ground-truth formalizations achieves good structure (ISR/PR = 0.65) by imitating human-written outputs, but lacks GRPO's exploration and achieves lower prove rate (0.20 vs.\ 0.30). Exec-only GRPO achieves high accuracy but degraded structure. Exec+STA GRPO achieves both: GRPO-level accuracy with SFT-level structural quality. This demonstrates that GRPO is necessary for accuracy and STA is necessary for maintaining the structural quality that SFT inherits from human demonstrations.

% -----------------------------------------------------------
\subsection{Scaling Analysis}
\label{sec:scaling}

\begin{table}[t]
\centering
\small
\begin{tabular}{llccc}
\toprule
\textbf{Model} & \textbf{Condition} & \textbf{PR} & \textbf{ISR/PR} & \textbf{Held-out PR} \\
\midrule
Qwen3-4B & exec-only & 0.28 & 0.35 & 0.24 \\
Qwen3-4B & exec+STA & 0.30 & 0.65 & 0.26 \\
Qwen3-8B & exec-only & 0.38 & 0.40 & 0.32 \\
Qwen3-8B & exec+STA & 0.40 & 0.68 & 0.34 \\
\bottomrule
\end{tabular}
\caption{Scale comparison (4B vs.\ 8B). The STA benefit is approximately scale-invariant: ISR/PR improvement of ${\sim}0.28$ at both scales. Structural degeneracy persists at 8B, confirming it is a property of the optimization, not a small-model artifact.}
\label{tab:scaling}
\end{table}

Table~\ref{tab:scaling} demonstrates that STA's benefits are scale-invariant. Scaling from 4B to 8B increases prove rate by ${\sim}0.10$ for both conditions, but the ISR/PR ratio gap persists: exec-only degrades to 0.40 (slightly better than 4B's 0.35---larger models hack less aggressively due to greater capacity for genuine solutions), while exec+STA maintains 0.68. The relative ISR/PR improvement (${\sim}0.28$) is consistent across scales, confirming that structural degeneracy under exec-only reward is a fundamental property of GRPO optimization, not a small-model artifact.

% -----------------------------------------------------------
\subsection{Hacking Characterization}
\label{sec:hacking}

\begin{table}[t]
\centering
\small
\begin{tabular}{lcccc}
\toprule
\textbf{Shortcut Type} & \textbf{Base} & \textbf{EO} & \textbf{E+S} & \textbf{SFT} \\
\midrule
Tautological premise & 8 & 18 & 3 & 2 \\
Vacuous quantifier & 7 & 12 & 4 & 3 \\
Redundant clause & 5 & 8 & 6 & 4 \\
Over-constrained & 3 & 5 & 3 & 2 \\
Other/subtle & 2 & 2 & 2 & 1 \\
\midrule
\textbf{Total hack rate (\%)} & 25 & 45 & 18 & 12 \\
\bottomrule
\end{tabular}
\caption{Hacking taxonomy distribution (\%, 200 outputs per condition, $\kappa = 0.78$). EO = exec-only, E+S = exec+STA. Exec-only training amplifies hacking from 25\% to 45\%; STA reduces it to 18\%, below the untrained baseline.}
\label{tab:hacking}
\end{table}

Table~\ref{tab:hacking} provides qualitative evidence complementing the quantitative ISR metrics. Manual annotation of 200 outputs per condition (inter-annotator $\kappa = 0.78$) reveals the composition of structural degeneracy.

Exec-only training amplifies every hacking category relative to the base model: tautological premises increase from 8\% to 18\%, vacuous quantifiers from 7\% to 12\%. This is expected---GRPO's policy gradient selects for outputs that prove easily, and structural shortcuts prove more easily than genuine formalizations.

STA-rewarded training suppresses the most egregious categories: tautological premises drop from 18\% to 3\% (STA assigns ISR = 0 to tautology components), vacuous quantifiers from 12\% to 4\%. Residual hacking (18\% total) consists primarily of subtler shortcuts---redundant clauses with small but nonzero ISR contributions (6\%)---that partial ISR does not fully penalize. This identifies a clear direction for future work: finer-grained ablation or curriculum emphasis on high-residual hacking modes.

% ============================================================
\section{Discussion}
\label{sec:discussion}

\paragraph{STA in the ORM-PRM Taxonomy.}
STA fills a previously empty cell: dense + automatic rewards derived from solver semantics rather than human annotation. ORMs provide sparse (outcome-level) automatic feedback; PRMs provide dense (step-level) feedback but require manual labels. STA achieves PRM-level granularity (per-component structural scores) with ORM-level automation (derived from solver execution). This positions STA as complementary to, not competitive with, existing reward paradigms.

\paragraph{When Does STA Help Most?}
Our results suggest STA's benefit scales with ISR discriminability---the degree to which components are independently ablatable. FOL (discriminability 0.72) shows the strongest effect; SMT-LIB (0.35) the weakest. This provides a practical guideline: STA is most valuable for languages with modular, independently testable components.

\paragraph{Limitations.}
Our experiments use a single base model family (Qwen3) and a limited set of formal languages. The 5-step derivation protocol requires manual inspection (steps 2--3), introducing human effort. ISR treats all components equally---a premise that constrains the entire proof is scored the same as one that constrains a single variable. Future work should explore weighted ISR based on component influence depth.

% ============================================================
\section{Conclusion}
\label{sec:conclusion}

We introduced Solver Trace Ablation (STA), a methodology for deriving dense structural rewards from solver execution semantics via the RIP framework from mutation testing. Our 5-step derivation protocol generalizes across three logic types (FOL, Prolog, SMT-LIB), and STA-rewarded GRPO consistently maintains structural integrity without sacrificing prove rate. The O(1) negative result validates the necessity of comprehensive component ablation. Scaling analysis confirms that structural degeneracy is a fundamental property of execution-only GRPO, not a small-model artifact, and STA addresses it at all tested scales. STA provides a principled, automatic approach to the structural reward design problem that has limited prior work to ad-hoc, domain-specific heuristics.

\begin{thebibliography}{9}

\bibitem[Mellgren et~al.(2025)]{mellgren2025prolog}
E.~Mellgren, A.~Bj{\"o}rklund, and J.~Nivre.
\newblock Structure-aware reinforcement learning for Prolog synthesis.
\newblock In \emph{Proc.\ ACL}, 2025.

\bibitem[EVOM(2026)]{evom2026}
Cross-solver GRPO transfer with execution-verified outcome models.
\newblock In \emph{Proc.\ ICLR}, 2026.

\bibitem[Polu and Sutskever(2020)]{polu2020gpt}
S.~Polu and I.~Sutskever.
\newblock Generative language modeling for automated theorem proving.
\newblock \emph{arXiv:2009.03393}, 2020.

\bibitem[Le et~al.(2022)]{le2022coderl}
H.~Le, Y.~Wang, A.~D.~Gotmare, S.~Savarese, and S.~Hoi.
\newblock CodeRL: Mastering code generation through pretrained models and deep reinforcement learning.
\newblock In \emph{Proc.\ NeurIPS}, 2022.

\bibitem[Lightman et~al.(2023)]{lightman2023verify}
H.~Lightman, V.~Kosaraju, Y.~Burda, H.~Edwards, B.~Baker, T.~Lee, J.~Leike, J.~Schulman, I.~Sutskever, and K.~Cobbe.
\newblock Let's verify step by step.
\newblock In \emph{Proc.\ ICLR}, 2024.

\bibitem[Jia and Harman(2011)]{jia2011analysis}
Y.~Jia and M.~Harman.
\newblock An analysis and survey of the development of mutation testing.
\newblock \emph{IEEE TSE}, 37(5):649--678, 2011.

\end{thebibliography}

\end{document}
